\section{\sys}
\label{s:mvscan}

\subsection{Overview}

\sys is an annotation-free static may-analysis. It builds a state-annotated interprocedural CFG; infers relation hypotheses; summarizes write coverage and read-to-sink influence; enumerates context-compatible writers and consumers; and aggregates the surviving evidence around a writer-centered native candidate. The output is for review, not a proved bug.

\subsection{Relation hypotheses}
\label{s:relation-discovery}

\paragraph{Control co-influence.}
For each modeled control sink, \sys collects abstract variables whose read-derived values may reach the condition. After eligibility filtering and normalization, two or more variables form an origin-specific hypothesis. Co-influence is useful because it identifies one decision made from several values; it is not independent evidence that those values must obey an invariant.

\paragraph{Return-derived relations.}
The second source is a return site in a function with no directly detected storage writes. This is not a transitive purity guarantee: the helper can call a mutating callee. One returned component must carry at least two variables; a tuple does not qualify merely because separate components expose one variable each. The analysis tracks components, but if a selected component has no recorded facts it falls back to tuple-wide information and can therefore mix components. The frozen environment option retains its historical ``multi-return'' name, while the paper uses the more accurate term \emph{return-derived relation}.

Scalar storage, complete keyed variables, and abstract external-state observations are eligible. Constants and immutables are excluded, and a precise keyed variable shadows its coarse mapping base. Every hypothesis retains an \emph{origin identity} (origin plus normalized members) and a member-only equivalence identity used for aggregation.

\subsection{Writer evidence}
\label{s:writer-analysis}

At writer site \(w\), the detector contextualizes each write description and independently matches it against relation members. After removing recognized syntactic self-copy patterns, the union of matches is \(\widehat{U}(w,\rho)\). A site survives when \(\varnothing\subset\widehat{U}\subset M\); the omitted set is \(M\setminus\widehat{U}\).

\(\widehat{U}\) is an abstract match set, not the semantic change set \(\Delta_M\), and its members need not be jointly realizable. A single \cc{m[k]} write can independently match \cc{m[0]} and \cc{m[1]} under different assumptions. The self-copy filter is also textual and lowercases identifiers, so Solidity names that differ only by case can be conflated.

Key matching preserves complete nested paths, substitutes formal arguments, and emits witness-local equality constraints when terms may coincide. The term language is deliberately restricted. State-variable references used as keys are treated as syntactically fixed even though their runtime values can change or coincide. Parameter canonicalization uses single-origin provenance; a phi merge carrying one parameter origin is not proof that the merged value equals that parameter on every branch. External-return template substitution is not uniform across all variable kinds, so receiver and argument translation is not claimed complete.

The frozen detector then applies a \emph{whole-root write-coverage heuristic}. It filters a writer when the root summary's syntactic targets cover every relation member. Coverage is not reconciliation: the generated sets can retain self-copies, do not establish appropriate value changes, and do not prove the relation holds before a sink. This heuristic can therefore lose recall.

\subsection{Consumer evidence}
\label{s:consumer-analysis}

Read origins propagate through assignments, phi operations, references, tuple handling, resolved calls, and returns. Summaries record which parameters and physical read events may influence:

\begin{itemize}[leftmargin=*,nosep]
  \item control operands, including branches and \cc{require}/\cc{assert};
  \item recorded value operands of modeled storage writes; and
  \item external effects.
\end{itemize}

Destination-key-only dependence is not generally recorded as a storage-write sink. An influence edge is behavioral evidence, not proof of a defect or changed concrete outcome.

Acceptance requires a location-specific read of an omitted variable under an owner-qualified root. Owner qualification remains mandatory even when the optional early root-context sink gate is disabled; the option controls an earlier filter, not the later multi-variable test. This is root sensitivity, not full callsite or invocation sensitivity.

\subsection{Candidates and structural buckets}
\label{s:candidate-construction}

Native candidates are keyed by writer root, writer site, member-set equivalence, matched set, and omitted set. Reader witnesses retain roots, abstract storage domains, parameter bindings, senders, sink sites, and witness-local key constraints. Candidate-level constraints are a union for search and display; constraints from different witnesses are alternatives, not one conjunction.

The context-instance count is a projection over roots, storage domains, senders, and block identifiers. Argument-binding tuples remain in reader-witness records but are omitted from that count. Reported call provenance is a representative root-to-writer function chain chosen by the implementation's parent selection. It is not guaranteed shortest, feasible, binding-exact, or a complete writer-to-reader execution trace.

For audit sampling, each native record maps to one structural payload containing dataset, subject, revision role, writer owner and block, relation members, matched and omitted members, transaction-context owners, sink kinds, origin sites, retained witness sites and their relation evidence, key constraints, and relation-shape tags. The native identifier, compilation-unit identity, and other details are discarded; identical payloads are deduplicated. A native record is not split into several buckets. Thus 3,732 \config{B0} native records become 2,868 structural buckets through many-to-one merging. \emph{Structural-bucket precision} is the fraction of audited buckets validated as MV-SI, not precision over native reports or distinct vulnerabilities.

\subsection{Approximation and cost drivers}

The analysis can over-approximate because a relation may be incidental, paths may not coexist, key constraints may be infeasible, or reconciliation may occur. It can under-approximate because relation origins, calls, keys, inheritance exposures, external observations, and whole-root coverage are incomplete abstractions.

The dominant costs are repeated summary propagation over IR, relation registration against access indexes, writer--reader Cartesian products, and contextual member matching. Call resolution and contexts are bounded, but fixed-point rounds, set sizes, and fan-out vary by program. We therefore report measured time and memory distributions instead of asserting an underspecified formal polynomial bound.
